{
  // Session controls how MCX runs the simulation and saves its results.
  "Session": {

    // Base name used for output files.
    // For example: cube60.mc2 and cube60.mch.
    "ID": "cube60",

    // Number of Monte Carlo photon packets to launch.
    // This is not necessarily the number of physical photons.
    "Photons": 1000000,

    // false: ignore refractive-index mismatch at material boundaries.
    // MCX will not apply Fresnel reflection or Snell refraction there.
    //
    // Note: refractive index n can still affect photon travel time.
    "DoMismatch": false,

    // true: let MCX automatically select the CUDA thread count,
    // block size, and related execution settings.
    "DoAutoThread": true,

    // true: save the volumetric simulation output.
    // The actual quantity saved is selected by OutputType below.
    "DoSaveVolume": true,

    // true: save records for photon packets captured by detectors.
    //
    // The name is slightly misleading: this activates detected-photon
    // output, while SaveDataMask chooses which fields are recorded.
    "DoPartialPath": true,

    // true: normalize the volumetric result to a unit source.
    //
    // This makes simulations with different photon-packet counts
    // directly comparable, apart from Monte Carlo noise.
    "DoNormalize": true,

    // false: do not save diffuse reflectance in background voxels
    // adjacent to the simulated object.
    "DoSaveRef": false,

    // false: do not save detected-photon exit positions and directions.
    //
    // If enabled, this adds the X and V fields to detector output.
    "DoSaveExit": false,

    // false: do not save each detected packet's random-number state.
    //
    // Saved seeds can later be used for photon replay calculations.
    "DoSaveSeed": false,

    // false: do not save momentum-transfer information.
    //
    // This option is primarily used for diffuse correlation spectroscopy
    // and other calculations involving momentum transfer.
    "DoDCS": false,

    // false: do not calculate the initial specular reflection that may
    // occur when a source enters the domain through a refractive boundary.
    "DoSpecular": false,

    // "0": no debugging output or photon-trajectory debugging.
    "DebugFlag": "0",

    // Selects the fields saved for each detected photon.
    //
    // This is a bit mask:
    //   1 = D = detector ID
    //   2 = S = scattering counts in each medium
    //   4 = P = partial path length in each medium
    //   8 = M = momentum transfer
    //  16 = X = exit position
    //  32 = V = exit direction
    //  64 = W = initial photon weight
    //
    // 5 = 1 + 4, so this saves:
    //   D: detector ID
    //   P: partial path lengths
    //
    // The more readable equivalent is "DP".
    "SaveDataMask": "5",

    // Save volumetric data in the legacy MCX .mc2 format.
    //
    // An .mc2 file is essentially a raw array of 32-bit floats and
    // does not contain dimensions or other metadata.
    "OutputFormat": "mc2",

    // "x" means time-resolved fluence rate.
    //
    // Case is ignored, so "x" and "X" mean the same thing.
    // This is volumetric fluence—not the detector packet count.
    "OutputType": "x",

    // Initial seed for MCX's random-number generator.
    //
    // Keeping the same seed helps make repeated runs reproducible,
    // especially when the GPU/thread arrangement is also unchanged.
    "RNGSeed": 1648335518
  },

  // Defines the simulated photon time window.
  "Forward": {

    // Simulation start time in seconds.
    "T0": 0,

    // Simulation end time in seconds.
    // This is approximately 5 nanoseconds.
    "T1": 4.99999996961265e-9,

    // Width of each time gate in seconds.
    // This is also approximately 5 nanoseconds.
    "Dt": 4.99999996961265e-9
  },

  // Because (T1 - T0) / Dt = 1, this configuration produces
  // one time gate covering approximately 0–5 ns.

  // Optode contains optical sources and detectors.
  "Optode": {

    // Definition of the photon source.
    "Source": {

      // A pencil source launches photons from one point in one direction.
      // It has no finite beam width.
      "Type": "pencil",

      // Source position in grid units: [x, y, z].
      "Pos": [
        29, // x coordinate
        0,  // y coordinate: the source is on the y=0 side
        29  // z coordinate
      ],

      // Initial direction: [vx, vy, vz, focal length].
      "Dir": [
        0, // no x component
        1, // points in the positive y direction
        0, // no z component
        0  // focal length = 0; collimated launch
      ],

      // First set of source-specific parameters.
      //
      // A pencil source does not use these parameters, so they are zero.
      "Param1": [
        0,
        0,
        0,
        0
      ],

      // Second set of source-specific parameters.
      //
      // Also unused by a pencil source.
      "Param2": [
        0,
        0,
        0,
        0
      ]
    },

    // List of spherical detector apertures.
    "Detector": [

      // Detector 1
      {
        // Detector-center position in grid units: [x, y, z].
        "Pos": [
          29, // x coordinate
          90, // y coordinate
          29  // z coordinate
        ],

        // Detector radius in grid units.
        // With the default 1 mm grid spacing, this is 1 mm.
        "R": 1
      },

      // Detector 2
      {
        "Pos": [
          29, // x coordinate
          45, // y coordinate
          59  // z coordinate, near the upper z boundary
        ],

        // Spherical detector radius in grid units.
        "R": 1
      }
    ]
  },

  // Shapes constructs the voxelized geometry directly from JSON.
  "Shapes": [

    {
      // Documentation-only name for this shape definition.
      // MCX does not use this value in the simulation.
      "Name": "cubic60"
    },

    {
      // Set the coordinate origin for all subsequent shape commands.
      "Origin": [
        0, // x origin
        0, // y origin
        0  // z origin
      ]
    },

    {
      // Create the background voxel grid.
      "Grid": {

        // Fill every voxel in this grid with material label 1.
        //
        // Label 1 corresponds to Domain.Media[1].
        "Tag": 1,

        // Grid dimensions in voxels: [Nx, Ny, Nz].
        "Size": [
          60,  // 60 voxels along x
          120, // 120 voxels along y
          60   // 60 voxels along z
        ]
      }
    }
  ],

  // Domain defines the grid coordinate convention and optical properties.
  "Domain": {

    // true means that the lower corner of the domain is [0, 0, 0].
    //
    // Source, detector, and shape coordinates therefore use a
    // zero-based spatial origin.
    "OriginType": true,

    // Domain dimensions in voxels: [Nx, Ny, Nz].
    //
    // This should agree with Shapes.Grid.Size.
    "Dim": [
      60,  // x dimension
      120, // y dimension
      60   // z dimension
    ],

    // Optical-property table.
    //
    // The array index corresponds to the voxel tag:
    //   Media[0] = tag 0/background
    //   Media[1] = tag 1
    //   Media[2] = tag 2
    //
    // No LengthUnit is specified, so the default is 1 mm per voxel.
    // Therefore mua and mus are in inverse millimeters, 1/mm.
    "Media": [

      // Media[0]: background material and space outside the domain.
      {
        // Absorption coefficient, μa = 0 mm^-1.
        // No photon weight is absorbed in this material.
        "mua": 0,

        // Scattering coefficient, μs.
        //
        // 1.19209e-7 mm^-1 is extremely close to zero, so this
        // material is effectively non-scattering.
        "mus": 1.19209289550781e-7,

        // Scattering anisotropy factor.
        //
        // g = 1 represents perfectly forward scattering, although it
        // has almost no effect here because mus is nearly zero.
        "g": 1,

        // Refractive index of the background material.
        "n": 1
      },

      // Media[1]: material assigned to every voxel by Grid.Tag = 1.
      {
        // Absorption coefficient, μa ≈ 0.005 mm^-1.
        //
        // The absorption-only characteristic distance is:
        // 1 / μa ≈ 200 mm.
        "mua": 0.00499999988824129,

        // Scattering coefficient, μs = 0.1 mm^-1.
        //
        // The average distance between scattering events is:
        // 1 / μs = 10 mm.
        "mus": 0.1,

        // Anisotropy factor, g ≈ 0.01.
        //
        // This is nearly isotropic scattering:
        // there is only a very small forward preference.
        //
        // It is not a scattering angle in radians.
        "g": 0.00999999977648258,

        // Refractive index of this material.
        //
        // Because DoMismatch is false, MCX will not apply Fresnel
        // reflection or Snell refraction at its boundaries.
        // The index can still affect photon propagation time.
        "n": 1.3700000047683716
      },

      // Media[2]: a second possible material.
      {
        // Absorption coefficient, μa ≈ 0.002 mm^-1.
        "mua": 0.0020000000949949,

        // Scattering coefficient, μs = 5 mm^-1.
        //
        // Mean distance between scattering events:
        // 1 / μs = 0.2 mm.
        "mus": 5,

        // Strongly forward-scattering medium.
        //
        // The reduced scattering coefficient is:
        // μs' = μs(1-g) = 5(1-0.9) ≈ 0.5 mm^-1.
        "g": 0.8999999761581421,

        // Refractive index of medium 2.
        "n": 1
      }
    ]
  }
}
